\documentclass[12pt]{article}
\usepackage{amsmath}
\usepackage{amsfonts}
\usepackage{amssymb}
\usepackage{graphicx}
\usepackage{enumitem}

\begin{document}
	
	\title{Examen Diagnóstico de Métodos Numéricos y tasas de interés}
	\date{}
	\maketitle
	
	Nombre: \underline{\hspace{6cm}} 
	\\
	\\
	Fecha: \underline{\hspace{6cm}}
	
	\section*{Parte A: Interpolación}
	
	\begin{enumerate}
		\item Mencione los métodos de interpolación que conozca.
		
		\item Explica en qué consiste el método de interpolación polinómica y menciona una ventaja y una desventaja de su uso.
		
		\item Utiliza el método de interpolación linear y \textit{log-linear} para encontrar el valor de $f(3.5)$ dada la siguiente tabla:
		
		\begin{tabular}{|c|c|}
			\hline
			$x$ & $f(x)$ \\
			\hline
			0 & 1 \\
			1 & 3 \\
			4 & 10 \\
			\hline
		\end{tabular}
	\end{enumerate}
	
	\section*{Parte B: Optimización}
	
	\begin{enumerate}[resume]
		\item Mencione los métodos de optimización que conozca para minimizar/maximizar funciones.
		
		\item Encuentra el valor mínimo de la función $f(x) = x^3 - 6x^2 + 9x + 2$ utilizando el método de la derivada.
	\end{enumerate}
	
	\section*{Parte C: Encontrar Ceros}
	
	\begin{enumerate}[resume]
		\item Define el método de la bisección y menciona una situación en la que este método podría ser preferible sobre otros métodos de encontrar ceros.
		
		\item Utiliza el método de la bisección para encontrar una aproximación de la solución de la ecuación $x^2 - 4x + 3 = 0$ en el intervalo $[0, 3]$.
		
		\item Describe brevemente el método de Newton-Raphson y menciona una condición bajo la cual el método puede converger lentamente o incluso diverger.
		
		\item Aplica el método de Newton-Raphson para encontrar la raíz positiva de la ecuación $\cos(x) - x = 0$. Utiliza $x_0 = 1$ como punto de partida y realiza dos iteraciones.
	\end{enumerate}
	
	\section*{Parte C: Conceptos de tasas de interés}
	\begin{enumerate}
		\item Mencione las convenciones de conteo de días que conozca.
		
		\item Mencione las convenciones de mercado que conozca para expresar a una tasa de interés.
		
		\item Suponga que desea valuar un bono con pagos de cupones semestrales, el plazo de dicho bono es 2 años, la tasa cupón es del 6\% anual pagadera cada seis meses y el nocional es \$1,000,000.00 mxn. Si la \textit{Yield To Maturity} del bono es del 5.3\% anual, lineal \textit{ACT/360}, cuál es el precio del bono?
		
		\item Mencione las tasas de interés de referencia que conozca.
		
		\item Defina qué es un Futuro de tasas de interés y un \textit{Forward Rate Agreement} explicando las diferencias entre ambos instrumentos.
		
		\item Defina qué es un \textit{Interest Rate Swap}, mencionando sus principales características.
	\end{enumerate}
\end{document}
